\documentclass{article}
\usepackage{hyperref}

\title{Planned Integrations for v2.0}
\date{\today}

\begin{document}

\maketitle
\section{Chapter 3 Revamp (Peng, Druv, and Qing)}

\begin{itemize}

    \item (Peng and Qing) Section 3.2 (Big noise pursue low-dim submanifolds) Characterization fo denoisers for one subspace and a mixture of subspaces. Integrating with the existing material (simulation etc.). Connection to PCA in Chapter 2 and pave way for Chapter 4 (CRATE) and Chapter 5 (RAE) transformer-like operators. 
     
    Connect this to the continuation story (for constrained optimization in Chapter 6) and tell a deeper story about diffusion.

    \item Section 3.3 (small Noise for low-dim submanifolds from samples)     Integrate content from Qing's works on memorization/generalization as well as Druv/Sam's. Lossy coding etc. 
    
    \item A new subsection of Section 3.3. Integrated characterization of `` multiple roles of noises'' ideas: a new narrative to unify the compression and denoising perspectives. Link with ongoing research into this problem.
    \begin{itemize}
        \item \url{https://github.com/Ma-Lab-Berkeley/deep-representation-learning-book/issues/5#issuecomment-3277109413}
    \end{itemize}
    \item \url{https://arxiv.org/abs/2509.20511} (related to continuation/our Appendix B?)
    
    \item Section 3.4 Transform the data distribution to a better feature distribution. (comments on contrastive learning DINO, VAE and LeJepa \url{https://arxiv.org/abs/2511.08544})
    Perform a more holistic integration of the compression and denoising stories by adding a latent diffusion story into the book. This can link the compression and denoising stories together at a deeper level rather than their current slightly awkward juxtaposition. Druv has ideas for how to do this based on convo with Kevin Murphy. This will need to be integrated with how to learn such representation in Chapter 5.
    We could feature the book cover diagram here---clearly stating its meaning/interpretation (links 3.2, 3.3, 3.4).

\end{itemize}

\section{Chapter 4: Causal architectures (Ziyang and Druv)}

\begin{itemize}
    \item This could be an example, or a larger component added where relevant
    \item Good way to enhance the sequence content of the book. We don't have much after Ch1
    \item It could be placed in Ch4 (used in Ch7 as a simple application to text/motion)
\end{itemize}

\section{Chapter 5 Revamping (Sam, Ziyang, and Yi)}

\begin{itemize}
    \item Section 5.1 Introduction: restate the objective of representation learning via rate reduction and how to learn the low-dim distribution of the representation (via compression or denoising) from Chapter 3. In addition to Chapter 4, now we ensure consistency or self-consistency via autoencoding or closed-loop transcription. 
    
    \item Section 5.2 After idealistic models PCA, Nonlinear PCA, and SAE, let us address how to learn a consistent representation for real general data distribution. First transform the data distribution to a better representation (MCR2-like objectives in Chapter 3). (Clarify the purpose of VAE as a comment or footnote.) SimDINO does this in an open-loop fashion (self-supervise via augmentation cropping). We can also learn an auto-encoding via White-box MAE (self-supervision via completion). Again the architecture is justified via the objective of the feature representation. 
    Mention that we can also do this in a closed-loop transcription fashion (for self-consistency) in  Section 5.3. 
    
    Then, use the diffusion and denoising (Chapter 3) to learn an autoencoding for the learned representation (say via the RAE method: \url{https://arxiv.org/abs/2510.11690}, high-dim diffusion with learned representations from an autoencoder works/scales angle. Justify the transformer-like architecture from the perspective of mixture of Gaussians (Chapter 3). 

    \item Section 5.3. Update the story based on the above revision for the earlier section. 
    \item Likely a new section in the future: Work on the new autoencoding-centric Cortex-like architecture (distributed, parallel, layer-wise optimization), given appropriate progress, would fit in very well here. This is ongoing research. It would also connect to the themes of Ch8 (distributed learning)

    % \begin{itemize}
    %     \item The general restructuring flow of the chapter could mirror Lecture 9 slides. 
    %     \item The idea is to think of what Peter did in the CPAL 2024 paper with Yann/Yi on Closed Loop as an "idealized" instantiation of the general autoencoding paradigm
    %     \item And for "practical" instantiations of this, think of Ziyang's SimDINO and Saining's RAE.
    %     \item Restructure the chapter around this kind of flow (remember that we start the chapter with a desideratum for autoencoding around rate reduction, but then we don't really follow it in the rest of the chapter). Then VAE is just an example rather than a focal point. (affects the VQ-VAE point below?)
    %     \item (Think about the standard practice: train the VAE separate from the diffusion model... agrees with what we are doing.)
    %     \item Describing VQ-VAE offers a way to more deeply connect with the compression/rate distortion/covering content of Chapter 3, as well as clarify a practically-ubiquitous methodology (which is maybe often understood or misused)
    %     \begin{itemize}
    %         \item Cool paper on this \url{https://arxiv.org/abs/2410.06424}
    %     \end{itemize}
    % \end{itemize}
\end{itemize}

\section{Chapter 6 Editing (Sam, Yi)}
\begin{itemize}
    \item Flesh out a further detailed vignette about how to solve inverse problems for image restoration (simple example, MRI reconstruction) using one of the simple methods, after\ DPS. We could also feature some of Liyue's work (or maybe code).
    \item We can probably condense the current synthetic DPS example (it's very long now, remember Kevin's feedback) and augment with this one.
    \item Replace the egomotion problem (section 6.3.3) and leave it for Chapter 7 as a practical application for real-world data. 
\end{itemize}

\section{Chapter 7 New Applications (Team)}
\begin{itemize}
    \item Text (new causal architectures? Druv and Ziyang).
    \item Conditioned human motion generation (Brent) Student can practice in class!
    \item Conditioned 2D image generation (RAE, Ziyang wu and Peter Tong) Mention Michelangelo as text-conditioned image generation. CLIP (image-text alignment, Ziyang)
    \item Image-conditioned 3D object generation (CUPID, Shenghua and Yi). 
    \item After all these modules ready, re-order and re-organize Chapter 7. 
\end{itemize}


\section{Other Matters}
\subsection{Refactoring and reorganization}

\begin{itemize}
    \item Do we want to keep thinking about the ``linear format'' of theory + applications, vs. a more integrated version?
    \begin{itemize}
        \item With the ``linear book'', it's hard to really justify a fully-interleaved approach
        \item But we want to perhaps have ``vignette'' form of applications in earlier chapters, and spin them out in detail in later chapter 7
    \end{itemize}
    \item Do we want to change the flow from Ch3 to Ch4?
    \begin{itemize}
        \item The focus on LDRs in Ch3 objectives section is a bit of an anachronism wrt the rest of the book (at least we need some polishing)
        \item Most of the text is very tuned to the previous journal version presentation. We might want to tune this a bit?
    \end{itemize}
    \item Do we want to integrate Ch3 and Ch6 more deeply? (Response to Kevin's comment)
\end{itemize}

\subsection{Figures replacement}

\begin{itemize}
    \item Replace many placeholder figures with more finalized versions
\end{itemize}

\subsection{Deeper code/data integration and open-sourcing}

\begin{itemize}
    \item In general, we should rely as much on code to explain the mathematical ideas as formulas and algorithm environments
    \begin{itemize}
        \item For pedagogy and clarity
        \item This will require lots of addition of new content throughout the book
    \end{itemize}
    \item A pipeline needs to be implemented to have code files in some source format, which can be rendered either into latex documents (listings) or into the website (syntax highlighting)
    \begin{itemize}
        \item Might want backlinks to the source as well, for web, in case where snippets (for figures?) are given
        \item Brent is working on this
    \end{itemize}
    \item Code used to generate figures (already in the Github) should be featured more prominently, in some cases.
\end{itemize}

\subsection{Slides and teaching material open-sourcing}



\end{document}